\chapter{Timing, Calibration and Localization}

White Rabbit or IEEE 1588 HA can distribute deterministic time and frequency over fibre, but timing synchronization is not automatically carrier-phase coherence at Ka. The phase uncertainty is:
\[
\sigma_\phi=2\pi f\sigma_t.
\]
Every sample must be linked to a common epoch and corrected for measured channel delay:
\[
t_{i,j,n}=T_{\mathrm{epoch}}+\frac{n}{f_{s,i,j}}+\delta_{i,j},\qquad t_i^\ast=t_i-\tau_i(f,T).
\]
For an illuminator $T$, target $x$ and receiver $R_i$, the bistatic excess path is:
\[
\Delta\rho_i=\lVert x-T\rVert+\lVert x-R_i\rVert-\lVert T-R_i\rVert=c\Delta\tau_i.
\]
AOA, TDOA and FDOA are fused with covariance, propagation models and geometric dilution of precision.

All faces and receiver chains within a radome share the node's central time base. Across radomes, that base is referenced to the network common epoch: GNSS and a frequency reference discipline the local clock, White Rabbit distributes time and frequency, atomic holdover preserves continuity, and calibration injection measures RF and digital delays. These techniques operate simultaneously and are checked continuously; each observation timestamp carries its calibration version and uncertainty. Echoes associated across several antennas of one radome and across other radomes can therefore be reconciled on the same time axis.

\section{GNSS timing reference}
Each station provisionally includes at least two external, surveyed GNSS antennas (\texttt{TIM-001}) with clear sky view, separated from the VHF/UHF Yagis and other high-power RF sources. A multi-constellation receiver produces UTC, 1PPS, frequency reference and ephemeris data. The proposed 1PPS and 10 MHz outputs (\texttt{TIM-002}) discipline a local atomic clock or high-stability oscillator; the clock then distributes deterministic timing through White Rabbit to every face and ADC.

The GNSS antenna phase centre, cable delay, receiver delay and distribution delay are measured and versioned. GNSS provides absolute epoch and long-term frequency discipline, while a local atomic holdover preserves continuity during loss of GNSS. End-to-end RF calibration remains mandatory:
\[
t_i^\ast=t_i-\tau_{\mathrm{GNSS},i}-\tau_{\mathrm{cable},i}-\tau_{\mathrm{RF},i}-\tau_{\mathrm{ADC},i}.
\]
Known station coordinates allow consistency checks between surveyed position, satellite geometry and measured network timing, but do not by themselves eliminate multipath, antenna phase-centre error or analog delay.

\begin{figure}[H]
\centering
\includegraphics[width=0.95\textwidth]{en/fig07_sincronizacao.png}
\caption{Timing and calibration architecture: the two left sources represent absolute GNSS reference and local holdover, the centre is the disciplined clock, the upper and lower right branches distribute timing to acquisition, and the final block represents measured RF-delay correction.}
\label{fig:fig07_sincronizacao}
\end{figure}

\begin{figure}[H]
\centering
\includegraphics[width=0.95\textwidth]{en/fig08_geometria_bistatica.png}
\caption{Bistatic/multistatic geometry: the orange point is the illuminator, the upper polygon is the target, blue points are receivers, orange and red paths are illumination and target-to-receiver propagation, and the dashed arc represents an equal-delay locus.}
\label{fig:fig08_geometria_bistatica}
\end{figure}
